\section{Discussion and Immediate Next Steps}

The current repository state is strong enough to support a focused method
manuscript, but not yet strong enough to skip the calibration evidence. The
core narrative is already clear:

\begin{enumerate}[leftmargin=1.5em]
  \item build a hierarchy,
  \item annotate it with edge and sibling significance tests,
  \item correct for multiplicity, and
  \item translate the annotated tree into cluster assignments.
\end{enumerate}

What remains is to show that these steps jointly behave well when used as an
end-to-end inferential pipeline.

For the next writing iteration, the most useful additions would be:

\begin{itemize}[leftmargin=1.5em]
  \item a diagram of the gate pipeline,
  \item one worked numerical example on a tiny binary tree,
  \item a calibration figure for null data,
  \item an ablation table over sibling-calibration methods, and
  \item a benchmark table tied to the exact current configuration reported in
  Table~\ref{tab:defaults}.
\end{itemize}

In other words, the manuscript environment is now ready for two parallel
streams of work: continued drafting of the prose and deliberate generation of
results that are specific to the method variant represented in the current
code path.
